% !TeX root = ../main.tex
% ---------------------------------------------------------------------------
% Rock-Support Interaction and Support System Design
% Scope extension added from the classical-textbook crosscheck (HIGH priority
% item: "expand support design beyond chart reading").
% Maps to: Hoek & Brown (1980) Ch. 8 "Underground excavation support design";
% Brady & Brown (2006) Ch. 11 "Rock support and reinforcement".
% Rationale: the book currently derives support from empirical charts
% (Chapter~\ref{ch:empirical-support}) and treats interaction only inside the
% CCM (Chapter~\ref{ch:ccm}). This chapter makes support a stiffness,
% capacity and timing problem, and opens Part V ahead of the component
% chapters.
% ---------------------------------------------------------------------------
\chapter{Rock--Support Interaction and Support System Design}
\label{ch:support-interaction}


\section{Support terminology and design variables}
\label{sec:support-terminology-and-design-variables}

\subsection{Terminology: support, reinforcement and retention}
\label{sec:support-terminology}
\scopenote{Support applied to the boundary, reinforcement acting within the
rock mass, and retention holding loose material in place. A complete system
needs all three, and specifications routinely confuse them.}

\subsection{The three variables that govern a support system}
\label{sec:support-three-variables}
\scopenote{Stiffness, capacity and installation timing. Chart-based design
selects a support class; it does not, by itself, tell you whether the system
will attract more load than it can carry.}

\section{Ground reaction and support characteristics}
\label{sec:ground-reaction-and-support-characteristics}

\subsection{The ground reaction curve as a design tool}
\label{sec:grc-as-tool}
\scopenote{Recapitulating the ground reaction curve of
Chapter~\ref{ch:ccm} in support-design terms: the ground sheds load as it
converges, and the support decides how much load it accepts.}

\subsection{Support characteristic curves for real systems}
\label{sec:scc-real-systems}
\scopenote{Deriving stiffness and maximum pressure for shotcrete shells, steel
sets, concrete linings, grouted bolts and segmental rings, rather than
assuming an idealised linear support.}

\section{Combined systems and installation timing}
\label{sec:combined-systems-and-installation-timing}

\subsection{Combined and sequential support}
\label{sec:combined-support}
\scopenote{Adding stiffnesses in parallel for simultaneously installed
elements, and handling elements installed at different convergences; the
common error of summing capacities that are never mobilised together.}

\subsection{Installation timing}
\label{sec:installation-timing}
\scopenote{Using the longitudinal deformation profile to convert distance
behind the face into convergence already released, and the consequences of
installing too early or too late.}

\section{Stiffness and yielding}
\label{sec:stiffness-and-yielding}

\subsection{Why stiffer is not always better}
\label{sec:stiffness-tradeoff}
\scopenote{A stiff support installed early attracts high load; a yielding
support installed at the same location survives. The design objective is
equilibrium at an acceptable convergence, not minimum convergence.}

\subsection{Yielding support and controlled deformation}
\label{sec:yielding-support}
\scopenote{Sliding sets, deformation slots and compressible elements as a
means of reaching equilibrium below the support capacity in squeezing ground.
Detail in Chapter~\ref{ch:squeezing-rockburst}.}

\section{Pre-reinforcement and block support}
\label{sec:pre-reinforcement-and-block-support}

\subsection{Pre-reinforcement and pre-support}
\label{sec:pre-reinforcement}
\scopenote{Reinforcement installed before the rock is disturbed retains the
intact rock mass strength; the classical argument for cable bolting ahead of
the face and for umbrella arches.}

\subsection{Support of blocks and wedges}
\label{sec:block-support}
\scopenote{Support design for the discrete blocks identified in
Chapter~\ref{ch:ground-type-design}: required bolt force, orientation and
factor of safety, treated separately from stress-driven support.}

\section{Support classes and reconciliation}
\label{sec:support-classes-and-reconciliation}

\subsection{Support classes and the ground-support system}
\label{sec:support-classes}
\scopenote{Building a set of standard support classes for a project, mapping
them to ground conditions, and the contractual and construction reasons for
keeping the number small.}

\subsection{Reconciling empirical, analytical and numerical designs}
\label{sec:reconciling-support}
\scopenote{The three routes give different answers because they represent
different mechanisms. A procedure for deciding which governs, extending
Chapter~\ref{ch:integrated-design}.}

\section{Quality and verification}
\label{sec:quality-and-verification}

\subsection{Installation quality and as-built performance}
\label{sec:support-quality}
\scopenote{Grout voids, under-thickness shotcrete, untensioned bolts and late
installation. Support that is designed correctly and installed poorly is the
most common cause of underperformance.}

\subsection{Worked example: selecting and verifying a support system}
\label{sec:support-worked-example}
\scopenote{Take a support class from the empirical charts, build its
characteristic curve, check equilibrium against the ground reaction curve at
the actual installation distance, and confirm the factor of safety.
Calculation sequence in Appendix~\ref{app:calculations}.}
